\section{第一次做 CEO：情绪价值、Ego 和棋手心态}

本章从产品回到创始人机制。明超平第一次做 CEO，既享受融资和关注，也觉得如履薄冰。AI 应用窗口很短，产品形态快速变化，团队需要稳定情绪，也需要创始人不断对抗 Ego。这里不是个人性格八卦，而是创业组织如何在极高不确定性中保持判断质量。

\begin{figure}[H]
\centering
\includegraphics[width=0.94\textwidth]{ego-countermeasure.png}
\caption{对抗 Ego：第一次做 CEO，要用第三方视角、情绪价值和棋手心态约束自我。自制概念图，依据 02:14:50--02:36:57 对谈内容整理。}
\end{figure}

\begin{knowledgebox}{读图：CEO 机制也是产品机制}
第三方视角来自辩论，情绪价值来自组织管理，如履薄冰来自风险意识，对抗 Ego 来自自我校准，棋手心态来自复盘和顺序判断。它们共同决定团队能不能持续做正确选择。
\end{knowledgebox}

\subsection{用户大于 Ego}

明超平提到团队里有一句话：用户大于 Ego。每个人都有自己的观点，也会用自己的方式解释数据；但真正的校准来自和用户聊、泡在 Discord 里、看用户怎么创作、怎么分享、怎么评论。对 AI 应用公司来说，用户反馈尤其重要，因为新场景没有成熟指标，很多时候只能从真实使用里建立直觉。

\begin{warningbox}{Ego 会伪装成战略判断}
创始人的自信有用，但也容易把个人审美、融资顺利和外部赞美误当成用户需求。用户大于 Ego 的意思不是没有观点，而是观点必须被用户场景持续校准。
\end{warningbox}

\subsection{棋手和对弈的人}

访谈结尾的“下棋”比喻很关键。明超平说，重要的不是复刻最终棋谱，而是理解每一步为什么下在那里。产品也是如此：今天看到微信、抖音、Instagram、TikTok 的成品，并不代表照抄形态就能成功。关键是顺序、优先级和时机：什么现在做，什么三年后做，什么绝对不做。

\begin{importantbox}{产品优先级就是下棋顺序}
第 50 步下错，可能就没有第 100 步。AI 应用创业也是这样：早期先做创作容器、分发、Agent 网络，还是先做企业服务、订阅、模型能力，都会改变后续路径。
\end{importantbox}

\subsection{本章小结}

第一次做 CEO 的难点不是只有融资和招人，而是持续校准自己。EP101 的后半段说明，Agent 应用创业需要产品判断、组织情绪、用户校准和对抗 Ego 同时在线。

